% エピグラフ共通レイアウト設定
\setlength{\epigraphwidth}{0.72\textwidth}
\setlength{\epigraphrule}{0.3pt}
\renewcommand{\epigraphsize}{\small}

% ------------------------------------------------------------------------------
% 1. 日本語訳付きエピグラフ（日本語版標準）
% \ChapterEpigraphWithTranslation{原文}{訳文}{出典}
% ------------------------------------------------------------------------------
\newcommand{\ChapterEpigraphWithTranslation}[3]{%
  \begingroup
  \renewcommand{\textflush}{flushepinormal}%
  \epigraph{%
    % --- 原文 (欧文イタリック) ---
    {\itshape\small #1}\par
    \vspace{0.45em}%
    % --- 訳文 (日本語ゴシック / 控えめなサイズ) ---
    {\sffamily\footnotesize\gtfamily #2}%
  }{%
    % --- 出典 ---
    \upshape\footnotesize --- #3%
  }%
  \endgroup
  \vspace{1.5em}%
}

% ------------------------------------------------------------------------------
% 2. 原文のみエピグラフ（英語版、または訳文不要時）
% \ChapterEpigraph{原文}{出典}
% ------------------------------------------------------------------------------
\newcommand{\ChapterEpigraph}[2]{%
  \begingroup
  \renewcommand{\textflush}{flushepinormal}%
  \epigraph{%
    {\itshape\small #1}%
  }{%
    \upshape\footnotesize --- #2%
  }%
  \endgroup
  \vspace{1.5em}%
}
